Kurs:Algebraische Kurven (Osnabrück 2012)/Arbeitsblatt 29/latex
Kurs:Algebraische Kurven (Osnabrück 2012)/Arbeitsblatt 29/latex


\setcounter{section}{29}


  \zwischenueberschrift{Aufwärmaufgaben}


\inputaufgabe

{}

{


Es sei $K$ ein
\definitionsverweis {algebraisch abgeschlossener Körper}{}{.}
Zeige, dass eine
\definitionsverweis {ebene projektive Kurve}{}{}
mit jeder projektiven Geraden in der projektiven Ebene einen nichtleeren Durchschnitt hat.


}

{} {}


\inputaufgabegibtloesung

{}

{


Es sei


\mavergleichskette

{\vergleichskette

{ K
}

{ = }{ \Z/( 5 )
}

{  }{
}

{    }{
}

{   }{
}

}
{}{}{,}
wir betrachten die beiden affinen ebenen algebraischen Kurven


\mathdisp {C = V \left(  X^2+Y^2-1  \right)  \text{ und } D = V \left(  X^3-2Y^2+3  \right)} { . }


\aufzaehlungvierabc{Bestimme den Durchschnitt $C \cap D$.
}{Bestimme die Punkte aus

\mathl{V_+ \left(  X^2+Y^2-Z^2  \right) \setminus V \left(  X^2+Y^2-1  \right)}{.}
}{Bestimme die Punkte aus

\mathl{V_+ \left(  X^3-2Y^2Z+3 Z^3  \right)  \setminus V \left(  X^3-2Y^2+3   \right)}{.}
}{Ist

\mathl{V_+ \left(  X^2+Y^2-Z^2  \right)}{} der
\definitionsverweis {projektive Abschluss}{}{}
von $V \left(  X^2+Y^2-1  \right)$?
}


}

{} {}


\inputaufgabegibtloesung

{}

{


Es sei $K$ ein Körper. Zeige, dass sämtliche lokale Ringe der projektiven Geraden ${\mathbb P}^{1}_{K}$ isomorph zueinander sind. Man gebe eine möglichst einfache Beschreibung dieses Ringes.


}

{} {}


\bild{ \begin{center}

\includegraphics[width=5.5cm]{ \bildeinlesungsvg {Lemniscate_of_Bernoulli} {svg} }

\end{center}

\bildtext {Die Lemniskate von Bernoulli} }


\bildlizenz { Lemniscate of Bernoulli.svg } {} {Zorgit} {Commons} {PD} {}


\inputaufgabe

{}

{


Bestimme für die durch

\mathl{V \left(  \left(  X^2+Y^2  \right)^2-X^2+Y^2  \right)}{} gegebene Lemniskate von Bernoulli die Singularitäten sowie die unendlich fernen Punkte in ${\mathbb P}^{2}_{ {\mathbb C} }$. Berechne in all diesen Punkten die Multiplizität und die Tangenten.


}

{} {}


\inputaufgabe

{}

{


Betrachte die durch die homogene Gleichung


\mavergleichskettedisp

{\vergleichskette

{ ZX^2
}

{ =} { Y^3
}

{ } {
}

{ } {
}

{ } {
}

}
{}{}{}
gegebene projektive Kurve


\mavergleichskette

{\vergleichskette

{C
}

{ \subset }{ {\mathbb P}^{2}_{ K }
}

{  }{
}

{    }{
}

{   }{
}

}
{}{}{}
über einem Körper $K$ der
\definitionsverweis {Charakteristik}{}{}
$0$.


a) Bestimme die singulären Punkte der Kurve.


b) Zeige, dass durch die Zuordnung


\mathdisp {\varphi: (S,T) \longmapsto \left( T^3  , \,  ST^2  , \,  S^3                 \right) =(X,Y,Z)} {  }


eine wohldefinierte Abbildung
\maabbdisp {\varphi} { {\mathbb P}^{1}_{}  } { {\mathbb P}^{2}_{}
} {}
gegeben ist.


c) Zeige, dass die Bildpunkte von $\varphi$ auf der Kurve $C$ liegen.


d) Welche Punkte in ${\mathbb P}^{1}_{}$ entsprechen den singulären Punkten der Kurve $C$.


}

{} {}


  \zwischenueberschrift{Aufgaben zum Abgeben}


\inputaufgabe

{}

{


Es seien

\mathl{m+1}{}
\definitionsverweis {homogene Polynome}{}{}


\mathl{F_0 , \ldots , F_m}{} in

\mathl{n+1}{} Variablen gegeben, die alle den gleichen  Grad $d$ besitzen. Zeige, dass es eine
\definitionsverweis {offene Menge}{}{}


\mavergleichskette

{\vergleichskette

{ U
}

{ \subseteq }{ {\mathbb P}^{ n }_{ K}
}

{  }{
}

{    }{
}

{   }{
}

}
{}{}{}
gibt, auf der die Polynome einen Morphismus


\mathdisp {{\mathbb P}^{ n }_{ K} \supseteq U \longrightarrow {\mathbb P}^{ m }_{ K}} {  }


definieren.


}

{} {}


\inputaufgabe

{}

{


Es sei


\mavergleichskette

{\vergleichskette

{ P
}

{ = }{ (a_0 , \ldots , a_n)
}

{ \in }{ {\mathbb P}^{ n }_{ K}
}

{    }{
}

{   }{
}

}
{}{}{}
ein Punkt im
\definitionsverweis {projektiven Raum}{}{.}
Zeige, dass die Projektion des ${\mathbb P}^{ n }_{ K}$ auf ${\mathbb P}^{n-1}_{K}$ mit Zentrum $P$ durch die Matrix


\mathdisp {} {  }


gegeben ist, also durch die Abbildung


\mathdisp {\begin{pmatrix}  x_0  \\ x_1 \\ \vdots\\ x_n   \end{pmatrix} \longmapsto  \begin{pmatrix}  x_0  \\ x_1 \\ \vdots\\ x_n   \end{pmatrix}} { . }


}

{} {}


\bild{ \begin{center}

\includegraphics[width=5.5cm]{ \bildeinlesungpng {Tschirnhausen_cubic} {png} }

\end{center}

\bildtext {Die Tschirnhausen Kubik} }


\bildlizenz { Tschirnhausen cubic.png } {} {Oleg Alexandrov} {Commons} {PD} {}


\inputaufgabe

{}

{


Bestimme für die durch

\mathl{V \left(  X^3+3X^2-Y^2  \right)}{} gegebene

\definitionswortenp{Tschirnhausen Kubik}{} die Singularitäten unter Berücksichtigung der unendlich fernen Punkte. Bestimme die Tangenten in den Singularitäten und in den unendlich fernen Punkten.


}

{} {}


\inputaufgabe

{}

{


Bestimme für das durch

\mathl{V \left(  X^3+Y^3-3XY  \right)}{} definierte
Kartesische Blatt
die unendlich fernen Punkte in ${\mathbb P}^{2}_{ {\mathbb C}}$ und berechne die Multiplizität und die Tangenten in diesen Punkten.


}

{} {}


\inputaufgabe

{}

{


Man gebe für die projektive Lemniskate von Bernoulli


\mavergleichskettedisp

{\vergleichskette

{ V_+ \left(  \left(  X^2+Y^2  \right)^2-Z^2X^2+Z^2Y^2  \right)
}

{ \subset} { {\mathbb P}^{2}_{ K }
}

{ } {
}

{ } {
}

{ } {
}

}
{}{}{}
einen surjektiven Morphismus auf eine projektive Quadrik an. Wie viele Punkte der Lemniskate werden dabei auf einen Punkt der Quadrik abgebildet?


}

{} {}


<< | Kurs:Algebraische Kurven (Osnabrück 2012) | >>


PDF-Version dieses Arbeitsblattes


 Zur Vorlesung
(PDF)
